\documentclass[11pt,a4paper]{article}
\usepackage{fontspec}
\setmainfont{DejaVu Sans}[Scale=0.90]
\usepackage{amsmath,amssymb}
\usepackage[margin=2.2cm]{geometry}
\usepackage{booktabs,array,tabularx}
\usepackage{enumitem}
\usepackage{titlesec}
\usepackage{parskip}
\titlelabel{\thetitle.\quad}
\titleformat{\section}{\normalfont\large\bfseries}{\thesection.}{0.6em}{}
\titleformat{\subsection}{\normalfont\normalsize\bfseries}{\thesubsection}{0.6em}{}
\titlespacing*{\section}{0pt}{15pt}{5pt}
\titlespacing*{\subsection}{0pt}{11pt}{4pt}
\setlist[itemize]{leftmargin=1.4em,itemsep=1pt,topsep=3pt}
% Preserve fonts/margins; permit a final line-breaking pass for long prose.
\setlength{\emergencystretch}{2em}
\begin{document}
\begin{center}
{\Large\bfseries Comparaison avec les travaux antérieurs}\\[4pt]
{\normalsize Fermeture, thermoréflectance, profils gradués et identifiabilité}\\[2pt]
{\small 16 septembre 2026 --- conventions du projet}
\end{center}

\section{Périmètre et état des sources}
Cette comparaison sépare ce qui est établi dans une source lue, ce qui est
recalculé ici et ce qui reste inaccessible. Les contrôles indépendants sont
des reproductions ou des vérifications, sans revendication de nouveauté.

\begin{center}
\begin{tabularx}{\textwidth}{@{}lXX@{}}
\toprule
Source & Accès obtenu & Portée de la comparaison\\
\midrule
Hennessy--Myers (2021) & Manuscrit accepté intégral, 14 pages &
Équations, conditions de bord, limite 1D, mesure pondérée\\
Beardo et al. (2020) & Article intégral, annexes incluses &
Transducteur, interfaces, différence avec le modèle du projet\\
Sendra et al. (2022) & Manuscrit intégral & Coefficient longitudinal et limite collective\\
Krapez (2016), IJHMT & Résumé et introduction de l'éditeur &
Structure PROFIDT ; pas de vérification ligne à ligne du texte intégral\\
Krapez (2016), JPCS & Notice HAL, extraits du texte déposé par l'auteur &
Article de conférence distinct ; téléchargement intégral indisponible\\
Camacho et al. (2025) & Corps de l'article déjà local &
Supplément toujours manquant, pas de reproduction de ses tables\\
\bottomrule
\end{tabularx}
\end{center}
Les deux articles Krapez 2016 ne sont pas interchangeables.
L'accès web au PDF HAL/IOP n'a pas fourni le document intégral.
Cette limite d'accès reste ouverte ; elle ne transforme pas une lecture
du résumé en validation des équations de l'article.

\section{Hennessy et Myers : raccord explicite}
Le manuscrit utilise GK historique ($\alpha=2$), un demi-espace isotrope,
un chauffage surfacique gaussien et une condition de glissement tangentielle.
Le transducteur métallique est omis. Sa décomposition du flux contient un mode
longitudinal et un mode transverse ; seul le premier porte une divergence,
mais les conditions de bord couplent leurs amplitudes. La température mesurée
est une moyenne spatiale complexe, puis on prend son argument.
Les équations (11)--(18) permettent un contrôle direct de la limite 1D.
Source : doi:\allowbreak 10.1007/\allowbreak 978-3-030-64272-3\_2.

\subsection{Conversion recalculée dans nos conventions}
Le manuscrit choisit $\exp(-i\omega t)$ ; le projet choisit $\exp(+i\omega t)$.
Les amplitudes sont donc conjuguées à paramètres réels.
Avec $k$ le nombre d'onde radial et $p=i\omega$,
\[
 \gamma_P^2=k^2+\frac{p(1+p\tau_R)}{a+3\ell^2p},
 \qquad
 \gamma_S^2=k^2+\frac{1+p\tau_R}{\ell^2}.
\]
Les racines ont une partie réelle positive.
Pour $k=0$, le mode transverse ne contribue pas au flux normal ;
la solution décroissante et $q=-\lambda_{\rm eff}T_z$ donnent
\[
 Z_{\rm 1D}=\frac1{\lambda_{\rm eff}\gamma_P}
 =\frac1{b\sqrt p}\sqrt{\frac{1+p\tau_R}{1+3\ell^2p/a}},
 \qquad
 \varphi=-\frac\pi4+\frac12\arctan(\omega\tau_R)
                  -\frac12\arctan(3\omega\ell^2/a).
\]
C'est exactement la réponse du code avec $\tau_\ell=3\ell^2/a$.
Le test numérique compare cette réduction à l'API du projet.

\subsection{Facteur trois dans le manuscrit accepté}
L'équation (18) contient $6\ell^2/L_F^2=3\omega\ell^2/a$.
La phrase suivante annonce une résonance à $\tau_R a/\ell^2=1$.
Avec \emph{ses propres} équations (2), (17) et (18), la condition exacte est
$\tau_R a/\ell^2=3$. Ce désaccord a été contrôlé sur le rendu de la page~9.
On retient la condition algébrique et on ne transpose pas la phrase isolée.
Cela ne constitue pas une vérification de la version finale éditée.

\section{Beardo et al. : ce qui manque à une expérience réelle}
L'article sur le silicium inclut un métal Fourier, un substrat hydrodynamique,
une source dans le métal et des faisceaux finis. L'annexe A impose continuité
du flux normal, saut de température avec corrections de non-équilibre, et
glissement du flux tangent. Les paramètres et comparaisons concernent le Si,
pas l'AlN. Son accord expérimental ne valide donc pas nos paramètres.
Source : Phys. Rev. B \textbf{101}, 075303 (2020),
doi:\allowbreak 10.1103/\allowbreak PhysRevB.101.075303.

Conséquence pour notre travail : un modèle 1D avec une simple résistance $R$
ne reproduit pas ce problème complet. Avant de tester un temps de relaxation,
il faut construire et contrôler au minimum le modèle Fourier du transducteur,
des deux interfaces et des rayons optiques. Une correction de conductivité
ne doit pas absorber automatiquement une erreur de spot ou de phase instrumentale.

\section{Krapez : propriétés transformées et problème inverse}
L'article IJHMT 99, 485--503 (2016),
doi:\allowbreak 10.1016/\allowbreak j.ijheatmasstransfer.2016.03.122,
présente une construction conjointe des profils d'effusivité solubles et des
champs thermiques par Liouville puis Darboux, avec des assemblages de profils.
Le texte intégral n'a pas été obtenu ; aucune formule numérotée ne lui est
attribuée ci-dessous. L'article JPCS 745, 032059 (2016),
doi:\allowbreak 10.1088/\allowbreak 1742-6596/745/3/032059,
traite des exemples de demi-espace et possède la notice HAL hal-01394669.

\subsection{Calcul indépendant : ce que la transformation conserve}
Pour Fourier gradué, $\xi=\int_0^z a^{-1/2}du$,
$b=\sqrt{\lambda C}$, $s=\sqrt b$, $\psi=sT$ donnent
\[
 \psi''=(p+V)\psi,\qquad s''=Vs,\qquad V=s''/s,\qquad
 q=s'\psi-s\psi'.
\]
Si $h''=(V+p_*)h$ et $r=h'/h$, poser
\[
 \widetilde V=V-2r',\quad
 \widetilde\psi=\psi'-r\psi,\quad
 \widetilde s=s'-rs .
\]
En dérivant et en utilisant $r'+r^2=V+p_*$, on retrouve
$\widetilde\psi''=(p+\widetilde V)\widetilde\psi$ et
$\widetilde s''=\widetilde V\widetilde s$.
Le profil thermique admissible exige $\widetilde s\ne0$ dans le domaine,
$\widetilde b=\widetilde s^2>0$ et des conditions de bord transformées cohérentes.
La construction d'une solution exacte n'est pas une preuve d'unicité inverse.

\subsection{Exemple de contrôle avec effusivité de volume finie}
Partons de $V=0$ et $h=\cosh(\xi/c+u_0)$, $c,u_0>0$.
Une normalisation de la propriété transformée donne
\[
 s=A\tanh(\xi/c+u_0),\quad
 V=-2c^{-2}\operatorname{sech}^2(\xi/c+u_0),\quad
 b_\infty=A^2 .
\]
Avec $\nu=\sqrt p$ et $t=\tanh u_0$, une solution décroissante fournit
\[
 Z(0,p)=\frac{\nu+c^{-1}t}
 {A^2t\nu(c^{-1}+t\nu)} .
\]
On vérifie $Z\sim1/(b_\infty\sqrt p)$ à basse fréquence et
$Z\sim1/(b_0\sqrt p)$ à haute fréquence, avec $b_0=A^2t^2$.
Le test indépendant intègre $T'=-q/b$, $q'=-pbT$ depuis une profondeur
suffisante et compare l'impédance à cette expression.
Cet exemple est un calcul de contrôle de la méthode, pas une extension
déjà intégrée dans l'API des couches graduées.

Enfin toute fonction positive $a(\xi)$ donne
$z(\xi)=\int_0^\xi\sqrt a\,d\xi$,
$\lambda=b\sqrt a$, $C=b/\sqrt a$.
Un même $b(\xi)$ admet donc des profils physiques différents.
Ni Darboux ni une optimisation réussie ne suppriment cette liberté sans données
indépendantes. Les notes 05--06 traitent déjà les commentaires et la réponse
de Mandelis ; aucune nouvelle conclusion d'unicité n'est ajoutée ici.

\section{Fermeture et comparaison avec les résultats déjà acquis}
Sendra et al. (2022), doi:\allowbreak 10.1103/\allowbreak PhysRevB.106.155301,
équation (21), obtient $1/3$ en approximation de Debye tridimensionnelle
et discute la différence avec le $2$ historique. Notre note~14 retrouve ce
coefficient par projection sans trace. Les expressions générales du manuscrit
font intervenir la dispersion et les collisions ; elles n'autorisent pas
à convertir quelques lectures de temps modaux en paramètres uniques d'AlN.

La résonance de Fourier était déjà identifiée chez Kovács (2018),
arXiv:1804.05225. Nous reproduisons la propriété sous nos conditions de bord.
Chez Camacho et al. (2025), la phase du demi-espace Cattaneo est un contrôle
du coefficient de flux. Sa sensibilité conditionnelle ne suffit pas à une
estimation jointe des paramètres. Le supplément absent empêche une reproduction
numérique exacte de tous ses exemples ; la nouvelle étude sera un scénario
autonome et intégralement spécifié.

\section{Décisions pour la suite}
\begin{enumerate}
\item Conserver le solveur 1D en $(b,\xi_1,\tau_R,\tau_\ell)$, dont les
relations sont inchangées ; annoncer la convention avant toute conversion cinétique.
\item Employer la fermeture conservatrice pour une interprétation issue du
modèle gris énergétique. Conserver la forme historique comme comparaison.
\item Construire l'analyse FDTR Fourier axisymétrique, puis examiner les
corrélations avec les interfaces, le transducteur et la calibration.
\item Ne pas annoncer une validation de GK 3D : les conditions de bord
non-Fourier correspondantes ne sont pas implémentées dans cette nouvelle base.
\item Pour achever la comparaison documentaire exhaustive, obtenir le texte
IJHMT 2016 et le supplément Camacho ; les limites sont localisées, pas masquées.
\end{enumerate}
\end{document}
